\section{Extensión opcional (extra desafío): LoRA Final}

Tras identificar los mejores hiperparámetros en la fase de optimización, se procede a la ejecución del \textbf{Desafío LoRA Final} mediante el comando \texttt{uv run nlp-train-lora}. Esta fase representa la implementación de "producción", donde se aplica un entrenamiento extendido para consolidar el conocimiento del modelo.

\subsection{Implementación de "Producción" con MLX}

A diferencia de los ensayos rápidos de Optuna, este proceso ya no utiliza valores fijos definidos manualmente, sino que \textbf{lee automáticamente los mejores hiperparámetros persistidos en las bases de datos SQLite de Optuna}. En concreto, para cada modelo se carga el estudio correspondiente y se extraen \texttt{study.best\_params}, reutilizando los valores ganadores de \texttt{learning\_rate} y \texttt{rank} sobre una configuración base de entrenamiento final.

De esta forma, \texttt{uv run nlp-train-lora} actúa como una segunda fase del pipeline experimental: Optuna explora y persiste; LoRA Final recupera el mejor resultado y lo entrena de forma más prolongada.

\begin{itemize}
    \item \textbf{Rank ($r$):} valor óptimo recuperado desde \texttt{results/optuna/optuna\_<model>.db}.
    \item \textbf{Scale:} ajustado en MLX como \texttt{scale = 2r}, manteniendo una relación estable entre capacidad y magnitud de la adaptación.
    \item \textbf{Learning Rate:} valor óptimo recuperado automáticamente del mejor \textit{trial}.
    \item \textbf{Capas Objetivo:} últimas capas lineales configuradas mediante LoRA en MLX.
    \item \textbf{Iteraciones:} 200 iteraciones por modelo en la fase final.
\end{itemize}

\subsection{Persistencia de Artefactos Finales}

La fase final genera dos tipos de artefactos:
\begin{itemize}
    \item \textbf{Pesos LoRA finales en \texttt{.adapters/}:} cada modelo guarda sus pesos en una carpeta independiente, por ejemplo \texttt{.adapters/Qwen3.5-2B/lora/final\_lora/} y \texttt{.adapters/Llama-3.2-1B-Instruct/lora/final\_lora/}.
    \item \textbf{Trazas de configuración en \texttt{results/lora/}:} se almacena una copia YAML de la configuración efectiva utilizada para cada entrenamiento final.
\end{itemize}

Este diseño mantiene separada la persistencia de pesos y la persistencia de metadatos experimentales, facilitando la reproducibilidad del entrenamiento y la auditoría posterior.

\subsection{Eficiencia y Resultados}

\begin{table}[H]
    \centering
    \begin{tabularx}{\textwidth}{|l|X|X|}
        \hline
        \textbf{Métrica}          & \textbf{Qwen 3.5 (2B)}  & \textbf{Llama 3.2 (1B)}    \\ \hline
        Parámetros Entrenables  & 5.60M (0.29\%)          & 3.20M (0.31\%)                \\ \hline
        Tiempo de Entrenamiento  & 10.5 min                & 8.4 min                      \\ \hline
        Estatus Final           & Completado              & Completado                    \\ \hline
    \end{tabularx}
    \caption{Resultados del proceso de Fine-Tuning LoRA.}
    \label{tab:lora_final}
\end{table}

\subsection{Conclusión del Desafío}

La implementación de LoRA ha demostrado ser la estrategia más eficaz para adaptar modelos de gran escala en entornos con recursos de memoria limitados (18GB), logrando una reducción de la \textit{Loss} significativa sin degradar el conocimiento base del modelo pre-entrenado.
